\section{Conclusion}
\label{sec:conclusion}

The paper includes a study on the optimization of the PicoRV32 RISC-V processor's performance. In this study, a Kogge-Stone adder and a PCPI-based Vedic Multiplier have been used.

The study's results have demonstrated that the adder's logic depth is reduced from 32 levels to 6 levels for a 32-bit RISC-V processor. In the context of the Vedic Multiplier, the latency of the multiplication process is reduced from 36 cycles to 3 cycles at a frequency of 100~MHz. This shows a 12.0$\times$ improvement in the processor's performance. In the context of the study's verification, the results have demonstrated 0 failures among 70,598 simulation vectors. Moreover, the study's formal checks have demonstrated full success.

The study's results have demonstrated that the integration of high-speed arithmetic cores into lightweight open-source RISC-V cores is a viable option. In the context of the study's design direction, the results have demonstrated the viability of the design direction for the development of edge-AI-capable embedded SoCs. Future work will be conducted through the same PCPI design direction, with the development of quantized AI-friendly instructions.
